\documentclass[12pt,a4paper]{article}
\usepackage[T1]{fontenc}
\usepackage[utf8]{inputenc}
\usepackage[polish]{babel}
\usepackage{lmodern}
\usepackage[a4paper,margin=2.2cm]{geometry}
\usepackage{hyperref}
\usepackage{enumitem}
\usepackage{booktabs}
\usepackage{xcolor}
\usepackage{titlesec}
\usepackage{fancyhdr}

\hypersetup{colorlinks=true,linkcolor=black,urlcolor=blue}
\pagestyle{fancy}
\fancyhf{}
\fancyhead[L]{\small PodpisECC --- instrukcja}
\fancyhead[R]{\small Adrian Gargula}
\fancyfoot[C]{\thepage}
\renewcommand{\headrulewidth}{0.4pt}

\title{\textbf{PodpisECC}\\[0.4em]
\large Instrukcja użytkowania, weryfikacji wyniku\\i dalszego rozwoju}
\author{System podpisu cyfrowego ECC\\ECDSA / secp256r1 / SHA-256 / OpenSSL}
\date{2026}

\begin{document}
\maketitle
\tableofcontents
\newpage

\section{Czy coś jeszcze trzeba?}

Na zaliczenie projektu \textbf{masz już komplet}:
\begin{itemize}[leftmargin=*]
  \item działającą aplikację webową (C++ + OpenSSL),
  \item generowanie kluczy, podpis, weryfikację, SHA-256,
  \item hasło do klucza prywatnego, limity uploadu, powiadomienia,
  \item pomiar ECC vs RSA-3072,
  \item skrypt testów automatycznych,
  \item dokumentację LaTeX/PDF w folderze \texttt{korostil}.
\end{itemize}

\textbf{Opcjonalnie} (niekonieczne, ale ładnie wygląda na obronie):
\begin{itemize}[leftmargin=*]
  \item 2--3 zrzuty ekranu UI w dokumentacji projektowej,
  \item krótka prezentacja / demo na żywo (happy-path + test negatywny),
  \item uzupełnienie wniosków po pomiarach z benchmarku.
\end{itemize}

Nie musisz dodawać PKI, HTTPS ani GUI desktopowego --- to było poza zakresem.

\section{Uruchomienie aplikacji}

\begin{enumerate}[leftmargin=*]
  \item Otwórz PowerShell.
  \item Przejdź do katalogu projektu:
\begin{verbatim}
cd "C:\Users\Adrian\Desktop\Podpis cyfrowy ecc"
\end{verbatim}
  \item Jeśli jeszcze nie budowałeś albo zmieniłeś kod:
\begin{verbatim}
cmake -S . -B build -G "MinGW Makefiles" ^
  -DCMAKE_BUILD_TYPE=Release ^
  -DOPENSSL_ROOT_DIR="C:/Strawberry/c"
cmake --build build
\end{verbatim}
  \item Uruchom serwer:
\begin{verbatim}
.\build\bin\podpis_ecc.exe
\end{verbatim}
  \item W przeglądarce otwórz: \url{http://127.0.0.1:8080/}
  \item Zatrzymanie: w oknie serwera \texttt{Ctrl+C} albo zamknięcie procesu.
\end{enumerate}

Inny port (np.\ gdy 8080 zajęty):
\begin{verbatim}
.\build\bin\podpis_ecc.exe 9090
\end{verbatim}
wtedy UI: \url{http://127.0.0.1:9090/}

\section{Jak używać --- krok po kroku}

\subsection{Scenariusz A --- podpis i weryfikacja poprawna (happy-path)}

\begin{enumerate}[leftmargin=*]
  \item W sekcji \textbf{1. Generuj klucze} wpisz hasło (zalecane, min.\ 6 znaków) dwa razy.
  \item Kliknij \textbf{Generuj parę kluczy}.
  \item Pobierz \texttt{private.pem} i \texttt{public.pem} (linki w wyniku).
  \item Przygotuj dowolny plik tekstowy, np.\ \texttt{umowa.txt}.
  \item W sekcji \textbf{2. Podpisz dokument}:
    \begin{itemize}
      \item upuść / wybierz dokument,
      \item upuść / wybierz \texttt{private.pem},
      \item wpisz hasło (jeśli klucz zaszyfrowany),
      \item opcjonalnie: \textbf{Podgląd SHA-256},
      \item kliknij \textbf{Podpisz i pobierz .sig}.
    \end{itemize}
  \item Zapisz pobrany \texttt{document.sig}.
  \item W sekcji \textbf{3. Zweryfikuj podpis} wgraj:
    \begin{itemize}
      \item ten sam dokument,
      \item \texttt{public.pem},
      \item \texttt{document.sig}.
    \end{itemize}
  \item Kliknij \textbf{Weryfikuj}.
\end{enumerate}

\textbf{Oczekiwany wynik:} toast i panel --- \textit{podpis PRAWIDŁOWY} (kolor zielony).

\subsection{Scenariusz B --- wykrycie zmiany dokumentu}

\begin{enumerate}[leftmargin=*]
  \item Po udanym podpisie zmień 1 znak w pliku dokumentu (lub wgraj inną wersję).
  \item Zweryfikuj \textbf{starym} podpisem i tym samym kluczem publicznym.
\end{enumerate}

\textbf{Oczekiwany wynik:} \textit{podpis NIEPRAWIDŁOWY} --- to dowód kontroli integralności.

\subsection{Scenariusz C --- niewłaściwy klucz}

\begin{enumerate}[leftmargin=*]
  \item Wygeneruj \textbf{drugą} parę kluczy.
  \item Dokument podpisany kluczem A zweryfikuj kluczem publicznym B.
\end{enumerate}

\textbf{Oczekiwany wynik:} \textit{podpis NIEPRAWIDŁOWY}.

\subsection{Scenariusz D --- błędy użytkownika (powiadomienia)}

Sprawdź, że aplikacja komunikuje konkretną przyczynę:
\begin{itemize}[leftmargin=*]
  \item brak pliku dokumentu / klucza / podpisu,
  \item klucz publiczny wgrany jako prywatny (i odwrotnie),
  \item złe hasło do zaszyfrowanego PEM,
  \item zaszyfrowany klucz bez podanego hasła,
  \item złe rozszerzenie pliku (\texttt{.pem}, \texttt{.sig}, \ldots).
\end{itemize}

\subsection{Scenariusz E --- pomiar wydajności}

W sekcji \textbf{Wydajność: ECC vs RSA-3072} ustaw 2--3 rundy i kliknij \textbf{Uruchom pomiar}.
Porównaj czasy keygen / sign / verify oraz rozmiar podpisu (ECC zwykle krótszy i szybszy w keygen).

\section{Jak weryfikować wynik (checklista na obronę)}

\begin{table}[h]
\centering
\begin{tabular}{@{}ll@{}}
\toprule
\textbf{Test} & \textbf{Oczekiwanie} \\
\midrule
Generowanie kluczy & dwa pliki PEM, opcjonalnie ENCRYPTED \\
Podpis & plik \texttt{.sig}, niepusty \\
Weryfikacja OK & \texttt{valid: true} / PRAWIDŁOWY \\
Zmiana dokumentu & NIEPRAWIDŁOWY \\
Zły klucz publiczny & NIEPRAWIDŁOWY \\
Złe hasło & komunikat błędu (HTTP 400) \\
Historia sesji & zdarzenia bez treści klucza prywatnego \\
\bottomrule
\end{tabular}
\end{table}

To wystarczy, żeby pokazać: \textbf{uwierzytelnienie}, \textbf{integralność}, \textbf{niezaprzeczalność} (przy tajności klucza prywatnego).

\section{Automatyczne testy}

Serwer musi być uruchomiony (\texttt{podpis\_ecc.exe} na porcie 8080).

W \textbf{drugim} oknie PowerShell:
\begin{verbatim}
cd "C:\Users\Adrian\Desktop\Podpis cyfrowy ecc"
powershell -ExecutionPolicy Bypass -File .\scripts\test_flow.ps1
\end{verbatim}

Skrypt sprawdza m.in.:
\begin{itemize}[leftmargin=*]
  \item dostępność UI,
  \item generowanie kluczy z hasłem,
  \item SHA-256,
  \item podpis + weryfikację pozytywną,
  \item weryfikację negatywną po zmianie dokumentu,
  \item odrzucenie złego hasła,
  \item endpoint benchmarku.
\end{itemize}

Sukces: komunikat \texttt{Wszystkie testy przeszły.} oraz linie \texttt{OK ...}.

\section{Co jeszcze można poprawić / rozwinąć}

\subsection{Szybkie ulepszenia (1--2 wieczory)}
\begin{itemize}[leftmargin=*]
  \item zrzuty ekranu do dokumentacji LaTeX,
  \item tabela wyników benchmarku wpisana do rozdziału wniosków,
  \item eksport historii sesji do pliku JSON/CSV (bez kluczy),
  \item ciemniejszy / jaśniejszy motyw UI.
\end{itemize}

\subsection{Rozszerzenia funkcjonalne}
\begin{itemize}[leftmargin=*]
  \item podpis wielu plików naraz (batch),
  \item weryfikacja ``przeciągnij folder'',
  \item zapis metadanymi: kto / kiedy (lokalny log, nie PKI),
  \item porównanie z ECDSA na \texttt{secp256k1} (Bitcoin) obok P-256,
  \item prosty tryb ``tylko weryfikacja'' dla drugiej osoby (bez generowania).
\end{itemize}

\subsection{Bezpieczeństwo (poziom zaawansowany)}
\begin{itemize}[leftmargin=*]
  \item magazyn klucza w systemie (DPAPI / Credential Manager),
  \item HTTPS lokalnie (opcjonalnie),
  \item HSM / token USB --- poza typowym zakresem przedmiotu.
\end{itemize}

\subsection{Czego nie trzeba robić na zaliczenie}
\begin{itemize}[leftmargin=*]
  \item pełne PKI / certyfikaty CA,
  \item kwalifikowany podpis elektroniczny,
  \item własna implementacja arytmetyki krzywej zamiast OpenSSL.
\end{itemize}

\section{Struktura projektu (skrót)}

\begin{verbatim}
Podpis cyfrowy ecc/
  src/           C++ (OpenSSL + serwer HTTP)
  web/           UI (HTML/CSS/JS)
  scripts/       test_flow.ps1
  docs/          ta instrukcja (PDF)
  build/bin/     podpis_ecc.exe
\end{verbatim}

Dokumentacja akademicka: folder \texttt{korostil} na Pulpicie (\texttt{build/dokument.pdf}).

\section{Szybka ściąga}

\begin{enumerate}[leftmargin=*]
  \item \texttt{.\textbackslash build\textbackslash bin\textbackslash podpis\_ecc.exe}
  \item Przeglądarka $\rightarrow$ \url{http://127.0.0.1:8080/}
  \item Generuj klucze $\rightarrow$ podpisz $\rightarrow$ weryfikuj OK
  \item Zmień dokument $\rightarrow$ weryfikuj FAIL
  \item \texttt{.\textbackslash scripts\textbackslash test\_flow.ps1} przy działającym serwerze
\end{enumerate}

\vspace{1.5em}
\noindent\textit{Koniec instrukcji. Powodzenia na zaliczeniu.}

\end{document}
